\documentclass{amsart}
% For dealing with references we use the comment environment
\usepackage{verbatim}
\newenvironment{reference}{\comment}{\endcomment}
%\newenvironment{reference}{}{}
\newenvironment{slogan}{\comment}{\endcomment}
\newenvironment{history}{\comment}{\endcomment}

% For commutative diagrams we use Xy-pic
\usepackage[all]{xy}

% We use 2cell for 2-commutative diagrams.
\xyoption{2cell}
\UseAllTwocells

% We use multicol for the list of chapters between chapters
\usepackage{multicol}

% This is generally recommended for better output
\usepackage{lmodern}
\usepackage[T1]{fontenc}

% For cross-file-references

% Package for hypertext links:
\usepackage{hyperref}

% For any local file, say "hello.tex" you want to link to please

% Theorem environments.
%
\theoremstyle{plain}
\newtheorem{theorem}[subsection]{Theorem}
\newtheorem{proposition}[subsection]{Proposition}
\newtheorem{lemma}[subsection]{Lemma}

\theoremstyle{definition}
\newtheorem{definition}[subsection]{Definition}
\newtheorem{example}[subsection]{Example}
\newtheorem{exercise}[subsection]{Exercise}
\newtheorem{situation}[subsection]{Situation}

\theoremstyle{remark}
\newtheorem{remark}[subsection]{Remark}
\newtheorem{remarks}[subsection]{Remarks}

\numberwithin{equation}{subsection}

% Macros
%
\def\lim{\mathop{\mathrm{lim}}\nolimits}
\def\colim{\mathop{\mathrm{colim}}\nolimits}
\def\Spec{\mathop{\mathrm{Spec}}}
\def\Hom{\mathop{\mathrm{Hom}}\nolimits}
\def\Ext{\mathop{\mathrm{Ext}}\nolimits}
\def\SheafHom{\mathop{\mathcal{H}\!\mathit{om}}\nolimits}
\def\SheafExt{\mathop{\mathcal{E}\!\mathit{xt}}\nolimits}
\def\Sch{\mathit{Sch}}
\def\Mor{\mathop{\mathrm{Mor}}\nolimits}
\def\Ob{\mathop{\mathrm{Ob}}\nolimits}
\def\Sh{\mathop{\mathit{Sh}}\nolimits}
\def\NL{\mathop{N\!L}\nolimits}
\def\CH{\mathop{\mathrm{CH}}\nolimits}
\def\proetale{{pro\text{-}\acute{e}tale}}
\def\etale{{\acute{e}tale}}
\def\QCoh{\mathit{QCoh}}
\def\Ker{\mathop{\mathrm{Ker}}}
\def\Im{\mathop{\mathrm{Im}}}
\def\Coker{\mathop{\mathrm{Coker}}}
\def\Coim{\mathop{\mathrm{Coim}}}

% Boxtimes
%
\DeclareMathSymbol{\boxtimes}{\mathbin}{AMSa}{"02}

%
% Macros for moduli stacks/spaces
%
\def\QCohstack{\mathcal{QC}\!\mathit{oh}}
\def\Cohstack{\mathcal{C}\!\mathit{oh}}
\def\Spacesstack{\mathcal{S}\!\mathit{paces}}
\def\Quotfunctor{\mathrm{Quot}}
\def\Hilbfunctor{\mathrm{Hilb}}
\def\Curvesstack{\mathcal{C}\!\mathit{urves}}
\def\Polarizedstack{\mathcal{P}\!\mathit{olarized}}
\def\Complexesstack{\mathcal{C}\!\mathit{omplexes}}
% \Pic is the operator that assigns to X its picard group, usage \Pic(X)
% \Picardstack_{X/B} denotes the Picard stack of X over B
% \Picardfunctor_{X/B} denotes the Picard functor of X over B
\def\Pic{\mathop{\mathrm{Pic}}\nolimits}
\def\Picardstack{\mathcal{P}\!\mathit{ic}}
\def\Picardfunctor{\mathrm{Pic}}
\def\Deformationcategory{\mathcal{D}\!\mathit{ef}}

\setlength{\emergencystretch}{3em}
\begin{document}
\title[Stacks Project, Tag 0F34]{562 --- Universal openness of finitely presented locally quasi-finite morphisms descends through affine inverse limits}
\maketitle
\noindent Copyright (C) 2005--2025 Johan de Jong. Stacks Project authors. Source: \href{https://stacks.math.columbia.edu/tag/0F34}{Tag 0F34}.

\noindent Editorial difficulty note (not author text). Difficulty H2: Étale localization and finite quasi-compact covers reduce the problem to a finite finitely presented algebra. The proof encodes universal openness by the image of a universal linear combination of module generators, then descends both that open image and its surjective parametrization to a finite stage..

\noindent Editorial provenance note (not author text). History: excerpt from revision \texttt{a0444\allowbreak 6e57e\allowbreak c1fbc\allowbreak 252a8\allowbreak 71afc\allowbreak ec775\allowbreak 2fb28\allowbreak 07b14}, more-morphisms.tex, lines 21877--21957. Reference presentation was adapted; original-author context and core support blocks were retained or added. The mathematical wording is unchanged. Licensed under GNU FDL 1.2 or later; no invariant sections or cover texts. See ../licenses/COPYING. Context blocks reproduce author definitions and supporting results; they are not additional counted problems.

\begin{lemma}
\label{lemma-descend-quasi-finite-universally-open}
Let $S = \lim S_i$ be a limit of a directed system of schemes
with affine transition morphisms.
Let $0 \in I$ and let $f_0 : X_0 \to Y_0$ be a morphism of schemes over $S_0$.
Assume $S_0$, $X_0$, $Y_0$ are quasi-compact and quasi-separated.
Let $f_i : X_i \to Y_i$ be the base change of $f_0$ to $S_i$ and
let $f : X \to Y$ be the base change of $f_0$ to $S$.
If
\begin{enumerate}
\item $f$ is locally quasi-finite and universally open, and
\item $f_0$ is locally of finite presentation,
\end{enumerate}
then there exists an $i \geq 0$ such that $f_i$ is locally quasi-finite
and universally open.
\end{lemma}

\begin{proof}
By Limits, Lemma \href{https://stacks.math.columbia.edu/tag/094M}{lemma descend quasi finite (Tag 094M)} after increasing
$0$ we may assume $f_0$ is locally quasi-finite. Let $x \in X$.
By \'etale localization of quasi-finite
morphisms we can find a diagram
$$
\xymatrix{
X \ar[d] & U \ar[l] \ar[d] \\
Y & V \ar[l]
}
$$
where $V \to Y$ is \'etale, $U \subset X_V$ is open, $U \to V$ is finite,
and $x$ is in the image of $U \to X$, see
Lemma \href{https://stacks.math.columbia.edu/tag/02LK}{lemma etale makes quasi finite finite at point (Tag 02LK)}.
After shrinking $V$ we may assume $V$ and $U$ are affine.
Since $X$ is quasi-compact, it follows, by taking a finite disjoint
union of such $V$ and $U$, that we can make a diagram as above
such that $U \to X$ is surjective. By
Limits, Lemmas \href{https://stacks.math.columbia.edu/tag/01ZM}{lemma descend finite presentation (Tag 01ZM)},
\href{https://stacks.math.columbia.edu/tag/01Z4}{lemma descend opens (Tag 01Z4)},
\href{https://stacks.math.columbia.edu/tag/07RR}{lemma descend surjective (Tag 07RR)},
\href{https://stacks.math.columbia.edu/tag/01ZO}{lemma descend finite finite presentation (Tag 01ZO)},
\href{https://stacks.math.columbia.edu/tag/07RP}{lemma descend etale (Tag 07RP)}, and
\href{https://stacks.math.columbia.edu/tag/01Z6}{lemma limit affine (Tag 01Z6)} after possibly increasing $0$
we may assume we have a diagram
$$
\xymatrix{
X_0 \ar[d] & U_0 \ar[l] \ar[d] \\
Y_0 & V_0 \ar[l]
}
$$
where $V_0$ is affine, $V_0 \to Y_0$ is \'etale,
$U_0 \subset (X_0)_{V_0}$ is open, $U_0 \to V_0$ is finite,
and $U_0 \to X_0$ is surjective. Since $V_i \to Y_i$ is \'etale
and hence universally open, follows that it suffices
to prove that $U_i \to V_i$ is universally open for large
enough $i$. This reduces us to the case discussed in the next
paragraph.

\medskip\noindent
Let $A = \colim A_i$ be a filtered colimit of rings.
Let $A_0 \to B_0$ be a ring map. Set $B = A \otimes_{A_0} B_0$ and
$B_i = A_i \otimes_{A_0} B_0$. Assume $A_0 \to B_0$ is finite,
of finite presentation, and $A \to B$ is universally open.
We have to show that $A_i \to B_i$ is universally open for $i$ large enough.
Pick $b_{0, 1}, \ldots, b_{0, d} \in B_0$ which generate $B_0$
as an $A_0$-module. Set $h_0 = \sum_{j = 1, \ldots, d} x_jb_{0, j}$
in $B_0[x_1, \ldots, x_d]$. Denote $h$, resp.\ $h_i$ the image of $h_0$
in $B[x_1, \ldots, x_d]$, resp.\ $B_i[x_1, \ldots, x_d]$.
The image $U$ of $D(h)$ in $\Spec(A[x_1, \ldots, x_d])$ is open
as $A \to B$ is universally open. Of course $U$ is quasi-compact
as the image of an affine scheme. For $i$ large enough there
is a quasi-compact open $U_i \subset \Spec(A_i[x_1, \ldots, x_d])$
whose inverse image in $\Spec(A[x_1, \ldots, x_d])$ is $U$, see
Limits, Lemma \href{https://stacks.math.columbia.edu/tag/01Z4}{lemma descend opens (Tag 01Z4)}.
After increasing $i$ we may assume that $D(h_i)$ maps
into $U_i$; this follows from the same lemma by considering the
pullback of $U_i$ in $D(h_i)$. Finally, for $i$ even larger
the morphism of schemes $D(h_i) \to U_i$ will be surjective
by an application of the already used
Limits, Lemma \href{https://stacks.math.columbia.edu/tag/07RR}{lemma descend surjective (Tag 07RR)}.
We conclude $A_i \to B_i$ is universally open by
Lemma \href{https://stacks.math.columbia.edu/tag/0F33}{lemma characterize universally open finite (Tag 0F33)}.
\end{proof}
\end{document}
